% Semantic Web Journal: A formal knowledge graph for the Rub' al-Khali
% desert ecosystem using the Semanticscience Integrated Ontology.
%
% This paper describes the KNOWLEDGE GRAPH ENGINEERING: ontology design
% patterns, OWL-Horst materialization, ELK validation, SPARQL competency
% questions, web portal, and design rationale. The underlying dataset
% is described in a companion GigaScience Data Note.
%
% Target: SWJ "Descriptions of Datasets" track or regular paper.

\documentclass[11pt,a4paper]{article}
\usepackage[utf8]{inputenc}
\usepackage[T1]{fontenc}
\usepackage[margin=1in]{geometry}
\usepackage{graphicx}
\usepackage{booktabs}
\usepackage{longtable}
\usepackage{hyperref}
\usepackage{gensymb}
\usepackage{amsmath,amssymb}
\usepackage{listings}
\usepackage{xcolor}
\usepackage{cite}

\graphicspath{{../}}

\title{A formal knowledge graph for integrating metagenomics and
  geochemistry of the Rub' al-Khali desert using the Semanticscience
  Integrated Ontology}

\author{\textit{Author list to be completed}}

\date{}

\begin{document}
\maketitle

\begin{abstract}
\noindent
Integrating multi-omic data with environmental metadata requires a
unified data model that can harmonize heterogeneous observation types,
enforce structural constraints at ingest, support logical consistency
checking, and provide semantic query answering at scale. We present a
formal knowledge graph that integrates 16S~rRNA gene amplicon
profiles, X-ray fluorescence measurements, field observations, and
climate records from a five-expedition sampling campaign across the
Rub' al-Khali desert (the companion data paper reports the expedition
and data acquisition in detail). The knowledge graph uses the
Semanticscience Integrated Ontology (SIO) as the upper-level
framework, with 302 domain classes grounded in ChEBI, ENVO, PATO, the
Units Ontology, and NCBI Taxonomy. We formalize eight reusable
ontology design patterns---for measurement, sampling, experimental
protocol, taxonomic abundance, X-ray fluorescence geochemistry,
sampling site with biome assignment, the sequencing workflow, and
sequencing quality control---in description logic and align them
explicitly with SIO's \textsc{Process--hasInput/hasOutput--Object},
\textsc{Collection--hasMember--Thing}, and
\textsc{Quantity--refersTo--Quality} idioms. A Shape Expressions
(ShEx) validation step enforces structural shapes on every RDF batch
before it is loaded into the triple store, and the ELK reasoner
confirms logical consistency of the full knowledge graph (zero
unsatisfiable classes) both before and after materialization. An
OWL-Horst (pD*) forward-chaining materialization workflow in Virtuoso
physically expands the graph from 37.0~million asserted triples to
117.9~million total triples (a 3.18$\times$ expansion) in
approximately 114~minutes. We evaluate the deployed resource through
nine SPARQL competency questions with sub-second to second-range
latencies, and we serve the knowledge graph through a public
SPARQL~1.1 endpoint and a web portal with an autocomplete-aware query
editor. We discuss the design rationale for every major choice,
connect each choice to the underlying challenge, and identify
limitations and directions for future work.

\vspace{0.5em}
\noindent\textbf{Keywords.} knowledge graph; SIO; OWL; ontology design
patterns; OWL-Horst; materialization; ELK reasoner; ShEx;
Virtuoso; SPARQL; desert microbiome; Rub' al-Khali
\end{abstract}

% ===================================================================
\section{Introduction}
% ===================================================================

Characterizing a complex ecosystem requires integrating heterogeneous
multi-scale data: microbial community profiles derived from DNA
sequencing, elemental soil composition from geochemical analyses,
spatiotemporal environmental variables, and provenance metadata that
links physical samples to their digital
representations~\cite{fierer2017}. Traditional approaches store these
data types in separate, disconnected repositories, which limits
cross-domain queries and reduces reproducibility. The emerging concept
of an ecosystem \emph{digital twin}~\cite{digitaltwin2021,
destinationearth,khan2025,iglesias2026} depends on precisely the kind
of unified, machine-interpretable, interoperable data model that is
missing for most natural environments; while biological systems and
specific regions have semantic digital
representations~\cite{biodigitaltwin}, no such resource exists for the
Rub' al-Khali, the largest contiguous sand desert on
Earth~\cite{mandaville_1986}, whose microbial communities and
biogeochemistry remain poorly characterized~\cite{makhalanyane2015}.

A companion data paper describes the underlying dataset: a
five-expedition soil sampling campaign (2023--2025) across a
1{,}043\,km transect of the Rub' al-Khali, yielding 16S~rRNA gene
amplicon profiles from 2{,}535 soil samples at 70 sampling sites,
X-ray fluorescence measurements for 91 chemical analytes from
249 measurement sessions, GPS-tagged environmental measurements
recorded at every site visit, and monthly and annual climate records
retrieved from the Open-Meteo database~\cite{openmeteo,
gigascience_companion}. This paper describes the \emph{semantic
layer}: the ontology engineering, the modular OWL/RDF architecture,
the reasoning workflow, and the evaluation of the deployed resource as
a Semantic Web application.

Constructing such a knowledge graph for an entire desert ecosystem
requires addressing several interlocking challenges that have shaped
the design of this work. First, the data are produced by heterogeneous
instruments and protocols (16S amplicon sequencing, X-ray fluorescence,
GPS-tagged field notes, gridded climate reanalysis), so a
\emph{unified data model} is needed that can represent samples,
processes, qualities, and quantities under a common upper-level
vocabulary rather than as disconnected tables. Second, taxonomic
labels are reported by multiple primer choices and classifier
workflows and must be \emph{harmonized} against a single reference
taxonomy without losing the original assertions. Third, because data
are deposited in batches over five field expeditions, \emph{schema
validation at ingest} is required to catch structural errors before
they propagate into the graph. Fourth, the knowledge graph must
support \emph{semantic query answering at scale}: with tens of
millions of triples, expressive OWL~2~DL reasoning is computationally
infeasible, so a tractable reasoning regime must be selected and its
trade-offs made explicit. Fifth, a \emph{user-facing query and
exploration interface} is required so that domain scientists who are
not fluent in SPARQL can interact with the resource. We return to
each of these challenges in Section~\ref{sec:discussion} and link
them to the corresponding design choices in the later sections.

This paper makes the following contributions to the Semantic Web
community. We present an OWL~2~DL knowledge graph that integrates a
desert ecosystem dataset under the Semanticscience Integrated
Ontology (SIO)~\cite{sio} as the upper-level framework, with domain
classes grounded in ChEBI~\cite{chebi}, ENVO~\cite{envo},
PATO~\cite{pato}, the Units Ontology~\cite{uo}, and NCBI
Taxonomy~\cite{ncbitaxon}. We formalize eight reusable ontology
design patterns in description logic and align them explicitly with
SIO's core relational idioms, accompanied by Turtle listings and
prose verbalizations so that readers can reproduce and extend
them. We integrate Shape Expressions (ShEx)~\cite{shex} validation
into an automated deployment workflow as a regression test against
structural drift. We implement OWL-Horst (pD*)~\cite{owlhorst}
forward-chaining materialization as three rounds of SPARQL Update
statements executed inside Virtuoso~\cite{virtuoso}, producing a
3.18$\times$ expansion from 37~million asserted to 118~million
total triples in approximately 114~minutes. We validate the logical
consistency of the full graph with the ELK reasoner~\cite{elk}
both before and after materialization and show that the
materialization introduces no disjointness violations. We report
compute-time metrics for the key deployment steps and per-query
latencies for a suite of nine SPARQL competency questions, and we
serve the resource through a public SPARQL~1.1 endpoint and a web
portal with an autocomplete-aware query editor. Finally, we discuss
the design rationale for every major choice, the current
limitations, and a sustainability plan for future data increments
and upstream ontology updates.

% ===================================================================
\section{Related work}\label{sec:related}
% ===================================================================

The Semantic Web community has produced a rich ecosystem of
standards, vocabularies, and tools that collectively cover parts of
the task of describing an ecological dataset. This section first
surveys those standards, then reviews the reasoning and deployment
literature that our work builds on, and finally articulates the gap
that motivates our design: no single existing standard covers the
full provenance chain from a physical soil sample to an integrated,
materialized knowledge graph that supports cross-domain query
answering, and combining them requires an integrative layer of the
kind we provide.

% -----------------------------------------------
\subsection{Upper-level ontologies for scientific data integration}
% -----------------------------------------------

The Semanticscience Integrated Ontology (SIO)~\cite{sio} is a
lightweight upper-level ontology that provides general-purpose classes
and relations applicable across scientific domains, including entities
(objects, processes, attributes), qualities, quantities, and
collections. SIO has been used to describe biomedical datasets,
laboratory experiments, and observational studies, and it is
compatible with several community vocabularies (OBO Foundry,
Units Ontology, ChEBI). We adopt SIO because its relational idioms
(\textsc{Process--hasInput/hasOutput--Object},
\textsc{Collection--hasMember--Thing}, and
\textsc{Quantity--refersTo--Quality}) generalize uniformly across the
sampling, laboratory, and computational processes that our dataset
requires. Recent work on simplified upper-level ontologies, notably
SULO~\cite{sulo2025}, takes this argument further: by strengthening a
few fundamental relations and reducing reliance on specialized
predicates, an upper-level ontology can lower the modelling effort
for new domains while preserving interoperability. Our design
patterns are compatible with this direction and could be re-grounded
in SULO with limited refactoring.

% -----------------------------------------------
\subsection{Related standards for provenance, observations, and
  sampling metadata}\label{subsec:related_standards}
% -----------------------------------------------

Several Semantic Web standards address subsets of the problem we
solve. We briefly review the most relevant, grouped by concern, and
explain in each case what the standard covers and what it leaves
uncovered.

\paragraph{Provenance: PROV-O.}
The W3C PROV-O~\cite{provo} recommendation defines a small number of
core classes (\texttt{Activity}, \texttt{Entity}, \texttt{Agent}) and
relations (\texttt{used}, \texttt{wasGeneratedBy},
\texttt{wasDerivedFrom}, \texttt{wasAssociatedWith}) that are
sufficient to record \emph{that} one data item was derived from
another, by whom, and when. PROV-O is deliberately domain-agnostic
and does not model what the activity \emph{does} beyond
``used input $X$, produced output $Y$.'' In our setting, every
scientific process---sampling, DNA extraction, library preparation,
sequencing, taxonomic classification, XRF analysis---is structurally
similar at the PROV-O level but semantically distinct: a sampling
process targets a site, an XRF process targets a sample and produces
typed concentration quantities, a bioinformatic workflow produces a
taxonomic dataset whose members are abundance observations. Encoding
these domain semantics in pure PROV-O forces the modeller to introduce
custom sub-classes of \texttt{Activity} and \texttt{Entity} for every
distinct process and data type. Our SIO-based patterns give this
structure directly in a single upper-level framework, and every
instance can be projected onto PROV-O as needed because SIO's
\textsc{Process--hasInput/hasOutput--Object} relations map
one-to-one onto PROV-O.

\paragraph{Observations and sensors: SOSA/SSN.}
The W3C/OGC SOSA (Sensor, Observation, Sample, Actuator)
vocabulary and the Semantic Sensor Network (SSN)
ontology~\cite{sosa} define classes for \texttt{Observation},
\texttt{Sensor}, \texttt{Sample}, \texttt{FeatureOfInterest}, and
\texttt{ObservableProperty}. SOSA/SSN is oriented primarily toward
Internet-of-Things deployments, continuous sensor streams, and the
relationship between a physical sensor and the property it measures.
It is well-suited for settings where the measurement process is a
thin wrapper around a sensor reading. Our dataset does not fit that
mould: we have one-off field measurements made with handheld
instruments, batch laboratory workflows, and computational analyses
whose ``observable property'' (for example, the relative abundance
of a taxon in a 16S~rRNA dataset) is itself the output of a multi-step
bioinformatics workflow. SOSA's \texttt{Sample} class is intentionally
thin and does not capture the soil-compartment / replicate structure
of our sampling design. Furthermore, SOSA does not model the
provenance chain from sample to digital data product. Our design
patterns could in principle be rephrased over SOSA by treating each
laboratory and computational step as a \texttt{hosts} relation on a
virtual sensor, but this mapping is unnatural and loses the
process-centric semantics that make the ODPs re-usable across
domains.

\paragraph{Experimental protocols: OBI.}
The Ontology for Biomedical Investigations (OBI)~\cite{obi} provides
a rich vocabulary for laboratory workflows, protocols, agents,
devices, and experimental designs in the biomedical
domain. OBI's coverage includes many of the process patterns we
need---material processing, library preparation, sequencing
assays---and its alignment with BFO and IAO gives it a sound
upper-level grounding. Two considerations kept us from adopting OBI
as our sole upper-level framework. First, OBI is heavily oriented
toward clinical and biomedical investigations; many of its classes
assume a human subject or a clinical assay context that does not
apply to environmental sampling. Second, OBI is large and has a
steep onboarding cost for a focused dataset like ours. We therefore
use OBI indirectly through SIO compatibility: several of our DNA
extraction and sequencing classes could be re-aligned with OBI
classes as a future extension, and we already import ChEBI and PATO
which are in the OBO family OBI uses.

\paragraph{Minimum metadata for sequencing: MIxS.}
The Genomic Standards Consortium's MIxS checklist~\cite{mixs}
specifies the minimum metadata fields required for sequence
submissions (MIMARKS for marker genes, MIGS/MIMS for genomes and
metagenomes). MIxS is a flat key-value checklist rather than an
ontology: it prescribes \emph{which} fields must be reported and
what controlled vocabularies to draw values from, but provides no
formal semantics and no inference. It is sufficient for ENA
submission but insufficient for cross-domain reasoning. A MIxS
record can be projected into our knowledge graph as a set of
triples about a sample, an environmental context, and an
associated sequencing run, but the inverse projection from our
knowledge graph to a MIxS record is trivial: every field the
checklist requires corresponds to a SPARQL query against an instance
of our sampling and sequencing patterns. In that sense we treat
MIxS as a downstream consumer of our knowledge graph rather than an
upstream data model.

\paragraph{Ecological observations: OBOE.}
The Extensible Observation Ontology (OBOE)~\cite{oboe} provides a
general structure for ecological observations based on four core
classes: \texttt{Observation}, \texttt{Entity}, \texttt{Measurement},
and \texttt{Characteristic}. OBOE has influenced the design of
SOSA/SSN and is conceptually close to the measurement pattern we
use. Two reasons kept us from adopting OBOE directly: its
community uptake has remained modest compared with SIO and SOSA,
and the observation-centric framing does not extend naturally to
the laboratory and computational processes that dominate our
provenance chain. Where an observation in OBOE is a single event,
our pipeline has many events---sampling, extraction, amplification,
sequencing, bioinformatic analysis---each of which generates its
own quantities and qualities. Modelling this as a chain of OBOE
observations is possible but loses the
\textsc{Process--hasInput/hasOutput--Object} abstraction that makes
the patterns composable.

\paragraph{Biodiversity records: Darwin Core.}
Darwin Core~\cite{darwincore} is the dominant community standard for
biodiversity occurrence records and underlies the Global
Biodiversity Information Facility (GBIF). Darwin Core is primarily
a flat record format with a controlled vocabulary (Darwin Core
Terms) rather than a semantic ontology, though a
Darwin-Core-in-RDF mapping exists. The scope of Darwin Core is the
occurrence record---``this organism was observed at this place at
this time''---and it does not model laboratory or computational
processes. It can be linked in as a complementary vocabulary for
the taxonomic component of our knowledge graph, and our per-sample
taxon observations could be projected into Darwin Core records for
GBIF deposit, but Darwin Core alone does not provide the integrative
data model we need.

\paragraph{Units of measurement: QUDT and OM.}
Both QUDT~\cite{qudt} and OM~\cite{om_ontology} provide comprehensive
RDF-based vocabularies for units of measurement, quantities, and
dimensions. They are strong alternatives to the Units Ontology
(UO)~\cite{uo}. We use UO because it is part of the OBO Foundry and
therefore integrates cleanly with ChEBI and PATO in a single
import closure. QUDT and OM could be swapped in without changing any
of our ontology design patterns; the measurement pattern
(Equation~\ref{eq:measurement}) is unit-vocabulary-agnostic.

\paragraph{Information artefacts: IAO.}
The Information Artifact Ontology (IAO)~\cite{iao} formalizes
information content entities, documents, data items, and
measurement data. IAO is widely used in OBO-family resources for
distinguishing the digital artefact that records a measurement
from the measurement process itself. Our knowledge graph makes this
distinction implicitly through the SIO \texttt{Quantity} class and
the data-property links back to literal values; a future alignment
of our quantity classes with IAO would make the distinction
explicit at the cost of an additional import.

% -----------------------------------------------
\subsection{What this work adds}\label{subsec:what_we_add}
% -----------------------------------------------

The preceding survey shows that each existing standard covers a
proper subset of the integration task. PROV-O provides provenance
relations but not process semantics. SOSA/SSN provides observations
but is oriented toward sensor streams and thin samples. OBI provides
biomedical laboratory workflows but is heavy and biomedical-focused.
MIxS provides a reporting checklist but no formal semantics. OBOE
provides ecological observations but has limited uptake and
observation-centric framing. Darwin Core provides biodiversity
records but no laboratory or computational provenance. None of these
standards alone provides a unified data model that can represent
sampling sites, field observations, laboratory processes, sequencing
workflows, bioinformatic analyses, and taxonomic abundance
observations under a common vocabulary, with full provenance, with
in-workflow structural validation, and with logical consistency
checking at deployment time. Four specific gaps drive our work:

\begin{itemize}

  \item \textbf{Heterogeneity of process types.} The ETL for a field
    expedition touches at least six distinct process categories:
    field sampling, field measurement, laboratory extraction,
    laboratory measurement, sequencing, and computational analysis.
    Each category is served by a different community standard in
    the literature (OBOE/SOSA for field, OBI for laboratory, MIxS
    for sequencing, PROV-O or OBOE for analysis). Combining them
    requires either a Frankenstein import (every standard contributes
    its own upper-level classes) or a thin integrative layer. We
    chose the latter: a single SIO-based upper layer under which each
    process category instantiates the same
    \textsc{Process--hasInput/hasOutput--Object} idiom. Our
    contribution is the specific set of eight design patterns
    (Section~\ref{subsec:odps}) that show how this layering works in
    practice, with formal DL and executable Turtle for every pattern.

  \item \textbf{Full provenance chain with uniform IRIs.} Existing
    resources that cover parts of this chain typically use different
    identifier schemes at different stages: the physical sample is
    identified by an institutional barcode, the DNA extract by a
    laboratory accession, the sequencing run by an ENA accession,
    the taxonomic classification by an OTU string, and the analysis
    output by a file path. Cross-referencing these requires manual
    joining. Our knowledge graph assigns every entity---site, visit,
    sample, DNA extract, amplicon library, sequencing run, FASTQ
    dataset, taxonomic observation, QC measurement---a single
    persistent IRI under the \texttt{https://rubalkhali.science/kb/}
    namespace, and connects them through SIO's
    \texttt{hasInput}/\texttt{hasOutput} relations. As a result, a
    single SPARQL traversal can walk from a taxon abundance back to
    the soil sample and on to the sampling site and the environmental
    measurements recorded at that visit, as demonstrated by
    Listing~\ref{lst:sparql_env_tax}. No existing standard provides
    this provenance chain out of the box; it has to be built by the
    integrator. We provide it.

  \item \textbf{Reasoning and validation as deployment steps.}
    None of the standards reviewed above specify how to validate
    data at ingest time or how to materialize entailments at
    deployment time. Community practice is to treat validation and
    reasoning as separate, optional, downstream tools. We make both
    first-class steps in the deployment workflow: every ABox module
    is validated against a ShEx shape before it reaches Virtuoso,
    and the deployed graph is materialized via an explicit OWL-Horst
    forward-chaining procedure that is scripted as SPARQL Update
    statements inside the triple store. The materialization choice
    reflects a specific engineering trade-off---ELK is used for
    consistency checking (sound and complete for OWL~2~EL
    classification), while OWL-Horst is used for forward-chaining
    ABox expansion (a different, incomplete, but practically
    sufficient semantics)---that existing standards do not prescribe.
    Our contribution is the concrete recipe (Section~\ref{subsec:owl_rdf})
    showing how to combine these two reasoners in a real Virtuoso
    deployment.

  \item \textbf{A reusable ontology design pattern library with
    verbalized listings.} Our eight design patterns are designed to
    be reused. Each is formalized in description logic, illustrated
    with a concrete Turtle listing, verbalized in prose, and
    documented in the deposited ontology with natural-language
    annotations. This combination---formal, instance-level, and
    prose---is uncommon in practice: community ontologies typically
    provide either formal axioms without examples or example files
    without formal patterns. The former is opaque to domain
    scientists; the latter is under-specified. We provide both, and
    we release the patterns as a standalone contribution so that
    authors of future environmental knowledge graphs can adopt them
    directly.

\end{itemize}

In short, SIO and the standards reviewed above give us the building
blocks: SIO for the upper-level relations, ChEBI and ENVO and NCBI
Taxonomy for the grounded domain terms, PROV-O for optional
projection, MIxS for optional downstream export. What we add is the
explicit, documented, reusable \emph{integration layer} that binds
them together for a real-world environmental dataset and runs as a
validated, materialized, queryable deployment.

% -----------------------------------------------
\subsection{Large-scale OWL reasoning}
% -----------------------------------------------

Reasoning over the combination of an expressive TBox and a
large ABox has been a recurring challenge in the Semantic Web
community for two decades. HermiT~\cite{hermit}, Pellet, and other
full OWL~2~DL reasoners implement tableau algorithms that are sound
and complete for the full profile but do not scale to the kind of
37-million-triple ABox we have here on commodity hardware. Three
broad strategies have been developed to address scale. The first is
\emph{materialization} of entailments under a restricted
semantics: ter~Horst's pD* / OWL-Horst rules~\cite{owlhorst} define a
set of RDF inference rules that cover most of RDFS, subproperties,
inverses, and transitivity, and can be applied forward-chaining to
populate an inferred graph. Virtuoso implements OWL-Horst as a
backward-chaining rule set that is activated at query time, but also
supports forward-chaining materialization via SPARQL Update. The
second strategy is \emph{query rewriting} under the OWL~2~QL profile,
which allows the reasoner to rewrite the user's query to account for
the TBox without materializing any new triples. The third strategy is
\emph{consequence-based reasoning} for restricted profiles such as
OWL~2~EL; the ELK reasoner~\cite{elk} is a polynomial-time
consequence-based classifier for OWL~2~EL that scales to very large
TBoxes with millions of named classes. In this work we combine two
of these strategies: we use ELK for logical consistency checking
and classification of the TBox (where its polynomial guarantees make
the check tractable), and we use OWL-Horst forward-chaining
materialization in Virtuoso to expand the ABox so that query
answering against the deployed endpoint can be served from physical
triples without paying a per-query reasoning cost. We report
quantitative results for both components below.

% -----------------------------------------------
\subsection{Triple store deployment}
% -----------------------------------------------

Several triple stores provide inference support and SPARQL query
answering at production scale. We deployed Virtuoso~\cite{virtuoso}
because it is available as free software under the GNU General Public
License, supports OWL-Horst forward-chaining materialization, and has
demonstrated operational maturity in large-scale linked-data
deployments.

% -----------------------------------------------
\subsection{Structural validation}
% -----------------------------------------------

Shape Expressions (ShEx)~\cite{shex} and the Shapes Constraint
Language (SHACL) are the two main W3C-compatible languages for
describing and validating RDF graph structures. We chose ShEx because
the data-generation workflow is batch-oriented and can invoke the
validator from the same Groovy/Jena toolchain that produces the
RDF. SHACL is an obvious alternative and is supported in-database by
several triple stores; a future port of the shapes from ShEx to
SHACL would enable in-database validation at load time. We return
to this point in the discussion.

% -----------------------------------------------
\subsection{Ecological and environmental knowledge graphs}
% -----------------------------------------------

Several groups have published linked-data resources for biodiversity,
biogeography, and environmental data (for example GBIF, GloBI, and
various OBI-aligned ecosystem studies). Compared with these
resources, our work differs in two respects. First, it targets a
single geographic region and a single coordinated sampling effort,
allowing the knowledge graph to encode the full provenance chain
from physical sample to digital data product with uniform IRIs
throughout. Second, the dataset tightly couples microbiome sequencing
with soil chemistry and environmental measurements, which is unusual
at this scale and supports cross-domain queries that are difficult
or impossible to express against the individual source
repositories.

% ===================================================================
\section{Dataset overview}
% ===================================================================

The companion data paper~\cite{gigascience_companion} describes the
expedition, sampling protocol, DNA extraction, sequencing workflow,
taxonomy harmonization, and environmental data retrieval in detail.
We summarize the relevant figures here so that this paper remains
self-contained.

The dataset is the product of five field expeditions conducted
between early 2023 and late 2025 along a 1{,}043\,km transect of the
Rub' al-Khali from the western escarpment (19.15\degree N,
45.17\degree E) to the eastern border with Oman (20.71\degree N,
54.73\degree E). At each of 70 sampling sites we collected bulk
surface, bulk subsurface (``deep''), and rhizosphere soil samples in
three biological replicates, for a total of 2{,}535 samples. We
generated 16S~rRNA amplicon sequencing profiles for every sample on
Illumina MiSeq, processed the reads with nf-core/Ampliseq +
DADA2~\cite{nfcore,ampliseq,dada2}, classified ASVs against the
Genome Taxonomy Database (GTDB)~\cite{gtdb}, and harmonized the
resulting classifications across primer/classifier combinations
against NCBI Taxonomy~\cite{ncbitaxon} and the iNaturalist taxonomy,
yielding 1{,}401{,}008 taxon--sample abundance observations. For a
subset of the sites we generated X-ray fluorescence measurements for
91 chemical analytes in 249 measurement sessions, and at every site
visit we recorded temperature, atmospheric pressure, and relative
humidity with handheld field instruments. We additionally retrieved
historical climate data from the Open-Meteo database for every site,
yielding 12{,}936 monthly and 582 annual climate records.

We distribute the integrated knowledge graph as nine modular OWL/RDF
files listed in Table~\ref{tab:modules}. The modules together encode
the full provenance chain from a physical sample to every digital
data product derived from it.

\begin{table}[h]
  \centering
  \caption{Distributed knowledge graph modules.}
  \label{tab:modules}
  \small
  \begin{tabular}{lp{0.58\textwidth}}
    \toprule
    File & Contents \\
    \midrule
    \texttt{rubalkhali.owl} & Base ontology (TBox); 302 domain
      classes grounded in SIO, ChEBI, ENVO, PATO, UO. \\
    \texttt{rubalkhali\_sites.owl} & 70 sampling sites with
      coordinates and ENVO biome annotations. \\
    \texttt{rubalkhali\_samples.owl} & 2{,}535 soil samples with
      compartment type and replicate identifier. \\
    \texttt{rubalkhali\_measurements.owl} & Environmental measurements
      and 12{,}936 monthly / 582 annual climate records. \\
    \texttt{rubalkhali\_xrf.owl} & 6{,}674 XRF measurement values
      across 91 analytes from 249 sessions. \\
    \texttt{rubalkhali\_dna.owl} & 2{,}410 DNA extracts with
      concentration measurements, agents, kits, protocols. \\
    \texttt{rubalkhali\_sra.owl} & 1{,}242 sequencing run records
      linked to ENA accessions. \\
    \texttt{rubalkhali\_qc.owl} & Sequencing quality metrics (GC
      content, duplicate rate, sequence counts). \\
    \texttt{rubalkhali\_taxonomy\_abox.owl} & 1{,}401{,}008
      taxon--sample abundance observations. \\
    \bottomrule
  \end{tabular}
\end{table}

% ===================================================================
\section{Knowledge representation}\label{sec:kr}
% ===================================================================

We developed a formal knowledge base to integrate the different data
types generated by the project. The semantic layer encodes the full
provenance chain from physical sampling to digital data products,
enabling cross-domain queries and ensuring data interoperability
through shared vocabularies.

% -----------------------------------------------
\subsection{Architecture}
% -----------------------------------------------

The knowledge base is constructed in description
logic~\cite{dlhandbook} using a modular architecture that strictly
separates the terminology (Terminological Box, TBox, or ontology) from
the observed facts (Assertion Box, ABox, or domain knowledge). The
TBox defines the ontology of the knowledge graph using
OWL~2~DL~\cite{owl2}, while the ABox contains the instance data
generated from the expedition. We deployed the complete knowledge
base across three named RDF graphs: an ontology graph (17{,}025
triples) containing the TBox, a knowledge base graph containing the
asserted instance data, and a taxonomy reference graph containing the
aligned taxonomy module.

% -----------------------------------------------
\subsection{Prefix conventions}\label{subsec:prefixes}
% -----------------------------------------------

For brevity, the Turtle and SPARQL listings throughout this paper use
the prefix declarations of Table~\ref{tab:prefixes}; these
declarations should be prepended to any listing that the reader
wishes to execute against the public endpoint. With these
declarations in place, identifiers such as
\texttt{https://rubalkhali.science/kb/RAK\_P100001} are abbreviated to
\texttt{rak:P100001} and
\texttt{http://semanticscience.org/resource/SIO\_000291} to
\texttt{sio:000291} without loss of information.

\begin{table}[h]
\centering
\caption{Prefix declarations used in all Turtle and SPARQL listings.}
\label{tab:prefixes}
\small
\begin{tabular}{ll}
\toprule
\textbf{Prefix} & \textbf{Namespace} \\
\midrule
\texttt{rak:}       & \texttt{https://rubalkhali.science/kb/RAK\_} \\
\texttt{sio:}       & \texttt{http://semanticscience.org/resource/SIO\_} \\
\texttt{pato:}      & \texttt{http://purl.obolibrary.org/obo/PATO\_} \\
\texttt{uo:}        & \texttt{http://purl.obolibrary.org/obo/UO\_} \\
\texttt{envo:}      & \texttt{http://purl.obolibrary.org/obo/ENVO\_} \\
\texttt{ncbitaxon:} & \texttt{http://purl.obolibrary.org/obo/NCBITaxon\_} \\
\texttt{obo:}       & \texttt{http://purl.obolibrary.org/obo/} \\
\texttt{rdf:}       & \texttt{http://www.w3.org/1999/02/22-rdf-syntax-ns\#} \\
\texttt{rdfs:}      & \texttt{http://www.w3.org/2000/01/rdf-schema\#} \\
\texttt{xsd:}       & \texttt{http://www.w3.org/2001/XMLSchema\#} \\
\texttt{geo:}       & \texttt{http://www.opengis.net/ont/geosparql\#} \\
\texttt{dcterms:}   & \texttt{http://purl.org/dc/terms/} \\
\bottomrule
\end{tabular}
\end{table}

% -----------------------------------------------
\subsection{Ontology engineering and design patterns}\label{subsec:odps}
% -----------------------------------------------

We use the Semanticscience Integrated Ontology (SIO)~\cite{sio} as
the upper-level framework. SIO provides general-purpose classes and
relations applicable across scientific domains, ensuring semantic
interoperability with other SIO-based datasets. We defined all
domain-specific classes as subclasses of SIO classes, and all
introduced relations are sub-relations of SIO relations. We mapped
classes to SIO and maintained full semantic mappings in our public
repository (Section~\ref{sec:availability}).

The core ontology comprises 302 distinct domain-specific classes
integrated directly under the SIO framework. To accurately represent
the complex relationships between environmental samples, laboratory
processes, and bioinformatic workflows, we incorporated 4 object
properties and 18 datatype properties. In total, the terminology
model encapsulates 1{,}290 formal axioms. This compact semantic
structure maintains precise alignment with our standardized schemas
while minimizing redundant inferences.

The following eight subsections present the ontology design patterns
used in the knowledge graph. Each pattern is formalized in description
logic, illustrated with a Turtle listing, and verbalized in prose
so that readers unfamiliar with OWL or Turtle can follow the
example. All patterns explicitly instantiate SIO's core relational
idioms (\textsc{Process--hasInput/hasOutput--Object},
\textsc{Collection--hasMember--Thing},
\textsc{Quantity--refersTo--Quality}), and later patterns reuse the
earlier ones as building blocks.

% ...........................
\subsubsection{Measurement pattern}

We model a quantitative measurement using SIO's ontology design
pattern for measurement processes: a measurement is a process that
generates a value for a specific quality of a target entity.

\begin{align}
    \label{eq:measurement}
    \text{MeasurementProcess} & \sqsubseteq \text{Process} \\
    \text{Quality} & \sqsubseteq \text{Attribute} \\
    \text{Quantity} & \sqsubseteq \text{InformationContentEntity} \\
    \text{MeasurementProcess} & \sqsubseteq \exists \text{hasTarget}.\text{Entity} \sqcap \exists \text{hasOutput}.\text{Quantity} \\
    \text{Quantity} & \sqsubseteq \exists \text{hasValue}.\text{xsd:double} \sqcap \exists \text{hasUnit}.\text{Unit}
\end{align}

We identify all units with classes from the Units Ontology (UO)~\cite{uo}.
Qualities are expressed using the Phenotype and Trait Ontology
(PATO)~\cite{pato} and linked to SIO as subclasses of the
\emph{Quality} class. Listing~\ref{lst:ttl_temp} illustrates an
environmental temperature measurement taken at a sampling site.

\begin{lstlisting}[caption={Turtle representation of an environmental
  temperature measurement.},
  label={lst:ttl_temp}, basicstyle=\ttfamily\scriptsize, frame=single]
rak:P000001 a rak:0000006 ; # Measurement Process
  sio:000291 rak:1000001 ; # hasTarget
  sio:000229 rak:4000001 . # hasOutput (Value)

rak:4000001 a rak:0000010 ; # Quantity
  sio:000221 uo:0000027 ; # hasUnit (degree Celsius)
  rak:2000003 "20.7"^^xsd:double ; # hasValue
  sio:000215 rak:5000001 . # refersTo (Quality)

rak:5000001 a pato:0000146 ; # Temperature Quality
  sio:000011 rak:1000001 . # isAttributeOf
\end{lstlisting}

In words, an instance \texttt{rak:P000001} of a measurement process
has as its target the sampling site \texttt{rak:1000001} and produces
as output the quantity \texttt{rak:4000001}. That quantity carries
the literal value \texttt{20.7}, is reported in degrees Celsius
(\texttt{uo:0000027}), and refers to a temperature quality
\texttt{rak:5000001} which is in turn an attribute of the same
site. Reading the chain in either direction therefore connects the
numeric reading to both the measured site and the abstract quality
that it quantifies.

% ...........................
\subsubsection{Sampling pattern}

A sampling event is a process that derives a sample from a feature of
interest (the site).

\begin{align}
    \label{eq:sampling}
    \text{SamplingProcess} & \sqsubseteq \text{Process} \\
    \text{Sample} & \sqsubseteq \text{MaterialEntity} \\
    \text{SamplingProcess} & \sqsubseteq \exists
                                 \text{hasTarget}.\text{Site} \sqcap
                                 \exists
                                 \text{hasOutput}.\text{Sample} \\
    \text{Sample} & \sqsubseteq \exists \text{isDerivedFrom}.\text{Site}
\end{align}

Listing~\ref{lst:ttl_sampling} illustrates the semantic representation
of a surface soil sample collected during Trip~1.

\begin{lstlisting}[caption={Turtle representation of a sampling process.},
  label={lst:ttl_sampling}, basicstyle=\ttfamily\scriptsize, frame=single]
rak:P100001 a rak:0000024 ; # Sampling Process
  sio:000291 rak:1000001 ; # hasTarget
  sio:000229 rak:7000001 . # hasOutput (Collection)

rak:6000001 a rak:0000021 ; # Deep Soil Sample
  sio:000252 rak:1000001 . # isDerivedFrom
\end{lstlisting}

Read in prose, an instance \texttt{rak:P100001} of a sampling process
targets the site \texttt{rak:1000001} and yields a collection
\texttt{rak:7000001} as its output; the deep soil sample
\texttt{rak:6000001} is in turn asserted to be derived from the same
site. The redundancy between \texttt{hasTarget} on the process and
\texttt{isDerivedFrom} on the sample is intentional: the former
records the action, the latter records the resulting provenance edge
directly on the material entity, so queries can traverse from sample
to site without joining through the sampling event.

% ...........................
\subsubsection{Experimental protocol pattern}

We model experimental protocols as processes specified by a protocol
entity. Each process (for example DNA extraction, library preparation)
consumes inputs and produces outputs, explicitly linking agents and
devices. The pattern instantiates SIO's general
\textsc{Process--hasInput/hasOutput--Object} idiom and adds a
specification link to a reusable protocol description.

\begin{align}
    \label{eq:protocol}
    \text{ProtocolExecution} & \sqsubseteq \text{Process} \\
    \text{Protocol} & \sqsubseteq \text{InformationContentEntity} \\
    \text{ProtocolExecution} & \sqsubseteq \exists
        \text{isSpecifiedBy}.\text{Protocol} \\
    \text{ProtocolExecution} & \sqsubseteq \exists
        \text{hasInput}.\text{MaterialEntity} \sqcap \exists
        \text{hasOutput}.\text{MaterialEntity}
\end{align}

Listing~\ref{lst:ttl_protocol} demonstrates the formalization of a
DNA extraction process.

\begin{lstlisting}[caption={Turtle representation of a DNA extraction protocol.},
  label={lst:ttl_protocol}, basicstyle=\ttfamily\scriptsize, frame=single]
rak:P200001 a sio:000994 ; # Scientific Protocol Execution
  sio:000339 rak:DNA_ext_powersoil ; # isSpecifiedBy
  sio:000230 rak:6000001 ; # hasInput (Soil Sample)
  sio:000229 rak:6800001 . # hasOutput (DNA Extract)

rak:6800001 a rak:0000040 ; # DNA Extract
  sio:000008 rak:4900001 . # hasAttribute (Concentration Quality)
\end{lstlisting}

In words, the protocol-execution instance \texttt{rak:P200001} is
specified by the named extraction protocol
\texttt{rak:DNA\_ext\_powersoil}, takes the soil sample
\texttt{rak:6000001} as input, and produces the DNA extract
\texttt{rak:6800001} as output; the resulting extract carries a
concentration quality (\texttt{rak:4900001}) which is itself
quantified through a separate measurement process following the
pattern of Equation~\ref{eq:measurement}.

% ...........................
\subsubsection{Taxonomic abundance pattern}

Taxonomic profiles represent collections of taxon observations
containing associated relative and absolute abundances. Instead of
modeling the feature table as a flat matrix, we formalize the entire
table as a bioinformatic workflow that produces a dataset containing
individual taxonomic observations. We use the merged taxonomy (see
the companion data paper~\cite{gigascience_companion}) to assign
specific NCBI or GTDB classes as the target of each observation.
For instance, the sequence analysis for a specific sample yields an
absolute sequence count associated with a taxonomic class from NCBI
Taxonomy. We construct an abundance quality that targets the explicit
taxonomic identifier, binding the observational count back to the
biological workflow. Formally, the pattern instantiates SIO's
\textsc{Collection--hasMember--Thing} and
\textsc{Quantity--refersTo--Quality} idioms:

\begin{align}
    \label{eq:taxabundance}
    \text{BioinformaticWorkflow} & \sqsubseteq \text{Process} \\
    \text{TaxonomicDataset} & \sqsubseteq \text{Collection} \\
    \text{BioinformaticWorkflow} & \sqsubseteq \exists
        \text{hasInput}.\text{SequenceDataset} \sqcap \exists
        \text{hasOutput}.\text{TaxonomicDataset} \\
    \text{TaxonomicDataset} & \sqsubseteq \forall
        \text{hasMember}.\text{AbundanceObservation} \\
    \text{AbundanceObservation} & \sqsubseteq \text{Quantity} \sqcap
        \exists \text{refersTo}.\text{AbundanceQuality} \\
    \text{AbundanceQuality} & \sqsubseteq \exists
        \text{isAttributeOf}.\text{Taxon}
\end{align}

\begin{lstlisting}[caption={Turtle representation of taxonomic
  abundance formalized from a feature table.},
  label={lst:ttl_tax}, basicstyle=\ttfamily\scriptsize, frame=single]
rak:P290001 a rak:0000071 ; # Bioinformatic Workflow
  sio:000230 rak:7900001 ; # hasInput (FASTQ Dataset)
  sio:000229 rak:7290001 . # hasOutput (Taxonomic Dataset)

rak:7290001 a rak:0000074 ; # Taxonomic Dataset
  sio:000059 rak:4400001 . # hasMember (Observation Value)

rak:4400001 a sio:000070 ; # Quantity (Absolute Count)
  sio:000215 rak:5400001 ; # refersTo
  rak:2000021 "15"^^xsd:integer . # absolute read count

rak:5400001 a rak:0000078 ; # Absolute Abundance Quality
  sio:000011 ncbitaxon:286 . # isAttributeOf (Pseudomonas)
\end{lstlisting}

Verbalized, the bioinformatic-workflow instance \texttt{rak:P290001}
consumes a FASTQ dataset (\texttt{rak:7900001}) and produces the
taxonomic dataset \texttt{rak:7290001}; that dataset has as one of its
members the quantity \texttt{rak:4400001}, which records an absolute
read count of \texttt{15} and refers to an absolute abundance quality
that is in turn an attribute of \emph{Pseudomonas}
(\texttt{ncbitaxon:286}). One row of a feature table is therefore
expanded into a four-node sub-graph that preserves both the classifier
output and its provenance.

% ...........................
\subsubsection{XRF geochemistry pattern}

We model XRF analysis as a process that targets a soil sample and
produces multiple measurement values, each corresponding to the
concentration of a specific chemical analyte. We map each analyte to
the Chemical Entities of Biological Interest (ChEBI)
ontology~\cite{chebi} and to the PubChem database~\cite{pubchem}. Out
of 93 monitored geochemical parameters, we mapped 80 analytes directly
to corresponding ChEBI classes and aligned 92 parameters to PubChem
Compound Identifiers. The XRF pattern is a specialization of the
measurement pattern (Equation~\ref{eq:measurement}), with the target
restricted to a soil sample and the produced quantity bound to an
analyte-specific quality:

\begin{align}
    \label{eq:xrf}
    \text{XRFMeasurementProcess} & \sqsubseteq \text{MeasurementProcess} \\
    \text{XRFMeasurementProcess} & \sqsubseteq \exists
        \text{hasTarget}.\text{SoilSample} \sqcap \exists
        \text{hasOutput}.\text{ConcentrationQuantity} \\
    \text{ConcentrationQuantity} & \sqsubseteq \text{Quantity} \sqcap
        \exists \text{refersTo}.\text{AnalyteQuality} \\
    \text{AnalyteQuality} & \sqsubseteq \exists
        \text{isAttributeOf}.\text{SoilSample}
\end{align}

\begin{lstlisting}[caption={Turtle representation of an XRF
  geochemistry measurement.},
  label={lst:ttl_xrf}, basicstyle=\ttfamily\scriptsize, frame=single]
rak:P300001 a rak:0000080 ; # XRF Measuring Process
  sio:000291 rak:6000001 ; # hasTarget
  sio:000229 rak:4300001 . # hasOutput

rak:4300001 a sio:000070 ; # Quantity
  sio:000215 rak:5300001 ; # refersTo
  sio:000221 uo:0000187 ; # Percent unit
  rak:2000012 "58.2"^^xsd:double . # Concentration

rak:5300001 a rak:0000167 ; # Analyte Quality Class
  sio:000011 rak:6000001 . # isAttributeOf
\end{lstlisting}

In words, the XRF measurement process \texttt{rak:P300001} targets
the soil sample \texttt{rak:6000001} and produces the quantity
\texttt{rak:4300001}, which records a concentration of
\texttt{58.2}~percent (\texttt{uo:0000187}) and refers to an
analyte-specific quality \texttt{rak:5300001}. Because the analyte
quality is typed with a dedicated subclass per element, the same
pattern instance can be reused unchanged across all 91 geochemical
analytes by varying only the quality class IRI.

% ...........................
\subsubsection{Sampling site and biome assignment pattern}

Each physical sampling location is represented as a persistent
individual that is grounded in the Environment Ontology (ENVO) by
linking to a biome and to one or more environmental features. We use
both an existential and a universal restriction so that the asserted
biome both holds and is the only admissible value, which preserves
the closed-world intent of the field record.

\begin{align}
    \label{eq:site}
    \text{SamplingSite} & \sqsubseteq \text{Site} \\
    \text{SamplingSite} & \sqsubseteq \exists
        \text{hasBiome}.\text{ENVO:Biome} \sqcap \forall
        \text{hasBiome}.\text{ENVO:Biome} \\
    \text{SamplingSite} & \sqsubseteq \exists
        \text{hasEnvironmentalFeature}.\text{ENVO:EnvironmentalFeature}
\end{align}

\begin{lstlisting}[caption={Turtle representation of a sampling site
  grounded in ENVO.},
  label={lst:ttl_site}, basicstyle=\ttfamily\scriptsize, frame=single]
rak:1000001 a rak:0000002 ; # Sampling Site
  rdfs:label "Site 11" ;
  geo:asWKT "POINT(50.123 22.456)"^^geo:wktLiteral ;
  rak:2000050 envo:01000219 ; # hasBiome (desert biome)
  rak:2000051 envo:00000086 . # hasEnvironmentalFeature (sand dune)
\end{lstlisting}

In words, the sampling site \texttt{rak:1000001} carries a label, a
WKT geometry, and is asserted both to belong to the ENVO desert biome
and to exhibit a sand-dune environmental feature. The combined
existential and universal restriction on \texttt{hasBiome}
(Equation~\ref{eq:site}) ensures that any reasoner treats the
asserted biome as the only admissible value for that site.

% ...........................
\subsubsection{Sequencing workflow pattern}

Sequencing data are produced by a chain of laboratory and
computational processes that the knowledge graph represents as a
sequence of \textsc{Process--hasInput/hasOutput--Object} steps: a DNA
extraction process consumes a soil sample and produces a DNA extract;
a library preparation process consumes the DNA extract and produces
an amplicon library; a sequencing process consumes the amplicon
library and produces a FASTQ dataset whose members are sequence reads.
Each process is specified by a protocol and may carry agents and
device participants, reusing the experimental protocol pattern of
Equation~\ref{eq:protocol}.

\begin{align}
    \label{eq:sequencing}
    \text{DNAExtractionProcess} & \sqsubseteq \text{ProtocolExecution} \\
    \text{DNAExtractionProcess} & \sqsubseteq \exists
        \text{hasInput}.\text{SoilSample} \sqcap \exists
        \text{hasOutput}.\text{DNAExtract} \\
    \text{LibraryPreparationProcess} & \sqsubseteq \text{ProtocolExecution} \\
    \text{LibraryPreparationProcess} & \sqsubseteq \exists
        \text{hasInput}.\text{DNAExtract} \sqcap \exists
        \text{hasOutput}.\text{AmpliconLibrary} \\
    \text{SequencingProcess} & \sqsubseteq \text{ProtocolExecution} \\
    \text{SequencingProcess} & \sqsubseteq \exists
        \text{hasInput}.\text{AmpliconLibrary} \sqcap \exists
        \text{hasOutput}.\text{FASTQDataset} \\
    \text{FASTQDataset} & \sqsubseteq \text{Dataset} \sqcap \forall
        \text{hasMember}.\text{SequenceRead}
\end{align}

\begin{lstlisting}[caption={Turtle representation of the sequencing
  workflow from soil sample to FASTQ dataset.},
  label={lst:ttl_sequencing}, basicstyle=\ttfamily\scriptsize, frame=single]
rak:P200001 a rak:0000044 ; # DNA Extraction Process
  sio:000339 rak:L000001 ; # isSpecifiedBy (PowerSoil protocol)
  sio:000230 rak:6000001 ; # hasInput (Soil Sample)
  sio:000229 rak:6800001 . # hasOutput (DNA Extract)

rak:P250001 a rak:0000045 ; # Library Preparation Process
  sio:000339 rak:L000010 ; # Illumina 16S Prep Guide
  sio:000230 rak:6800001 ; # hasInput (DNA Extract)
  sio:000229 rak:6900001 . # hasOutput (Amplicon Library)

rak:P280001 a rak:0000046 ; # Sequencing Process
  sio:000230 rak:6900001 ; # hasInput (Amplicon Library)
  sio:000229 rak:7900001 . # hasOutput (FASTQ Dataset)

rak:7900001 a rak:0000064 ; # FASTQ Dataset
  rdfs:seeAlso <https://www.ebi.ac.uk/ena/browser/view/ERR1234567> .
\end{lstlisting}

Read as a chain, the DNA extraction process \texttt{rak:P200001}
consumes the soil sample \texttt{rak:6000001} and produces the DNA
extract \texttt{rak:6800001}; that extract is the input of the library
preparation \texttt{rak:P250001}, which yields the amplicon library
\texttt{rak:6900001}; finally the sequencing process
\texttt{rak:P280001} consumes the library and produces the FASTQ
dataset \texttt{rak:7900001}, which is cross-referenced to its ENA
accession via \texttt{rdfs:seeAlso}. Each link in the chain is the
same \texttt{hasInput}/\texttt{hasOutput} pattern, so a single SPARQL
traversal can recover the full provenance from any node.

% ...........................
\subsubsection{Sequencing quality control pattern}

Sequencing quality control is represented as an additional
computational process that takes a FASTQ dataset as input and
produces multiple measurement values, each instantiating the
generic measurement pattern of Equation~\ref{eq:measurement}.
Metrics (sequence count, GC content, duplicate rate, read length)
are split into forward (R1) and reverse (R2) subclasses so that the
provenance of paired-end runs is preserved.

\begin{align}
    \label{eq:qc}
    \text{SequencingQCProcess} & \sqsubseteq \text{Process} \\
    \text{SequencingQCProcess} & \sqsubseteq \exists
        \text{hasInput}.\text{FASTQDataset} \sqcap \exists
        \text{hasOutput}.\text{QCMeasurementValue} \\
    \text{QCMeasurementValue} & \sqsubseteq \text{Quantity} \sqcap
        \exists \text{refersTo}.\text{QCQuality} \\
    \text{QCQuality} & \sqsubseteq \exists
        \text{isAttributeOf}.\text{FASTQDataset}
\end{align}

\begin{lstlisting}[caption={Turtle representation of a sequencing QC
  measurement (forward-read GC content).},
  label={lst:ttl_qc}, basicstyle=\ttfamily\scriptsize, frame=single]
rak:P350001 a rak:0000047 ; # Sequencing QC Process
  sio:000230 rak:7900001 ; # hasInput (FASTQ Dataset)
  sio:000229 rak:4500001 . # hasOutput (QC Value)

rak:4500001 a rak:0000086 ; # GC Content Value
  sio:000215 rak:5500001 ; # refersTo
  sio:000221 uo:0000187 ; # Percent
  rak:2000073 "54.2"^^xsd:double . # has GC content

rak:5500001 a rak:0000074 ; # GC Content Quality (Forward)
  sio:000011 rak:7900001 . # isAttributeOf (FASTQ Dataset)
\end{lstlisting}

In words, the QC process \texttt{rak:P350001} takes the FASTQ dataset
\texttt{rak:7900001} as input and produces the measurement value
\texttt{rak:4500001}, which records a GC content of
\texttt{54.2}~percent and refers to a forward-read GC-content quality
\texttt{rak:5500001} that is in turn an attribute of the same FASTQ
dataset. Splitting the quality class into forward and reverse
subclasses preserves the read direction, so that paired-end QC
metrics can be queried independently per orientation.

% -----------------------------------------------
\subsection{OWL and RDF implementation}\label{subsec:owl_rdf}
% -----------------------------------------------

We generated the OWL/RDF implementation using custom Groovy scripts
built on the OWL API~5~\cite{owlapi}. The knowledge base employs a
modular design, separating the base ontology, sites, samples,
measurements, XRF data, DNA extraction records, sequencing metadata,
quality control metrics, and taxonomic abundances into individual
ABox modules~\cite{hitzler_modularity}. Each module is generated by
an independent script, validated against its ShEx shape, and loaded
into Virtuoso as a separate transaction, so that a single data-type
update can be rolled out without rebuilding the full graph.

We use the OpenLink Virtuoso triple store~\cite{virtuoso} for data
hosting and query answering because it is available as free software
under the GNU General Public License, supports OWL-Horst
forward-chaining materialization, and has demonstrated operational
maturity in large-scale linked-data deployments. Virtuoso features an
inference engine that can be configured with a custom rule set; we
bind the project ontology graph as a rule set named
\texttt{rubalkhali} that covers all three named graphs (knowledge
base, ontology, taxonomy) and provides backward-chaining expansion
for rdfs:subClassOf, rdfs:subPropertyOf, owl:inverseOf,
owl:equivalentClass, and related OWL-Horst rules.

Prior to OWL-Horst inference, the base knowledge graph comprised
37{,}043{,}525 asserted triples across three named graphs (knowledge
base 36{,}978{,}835; taxonomy 42{,}406; ontology 16{,}643), with
taxonomy assertions accounting for approximately 36.5~million of the
knowledge base graph. During deployment, the workflow applies
OWL-Horst (pD*) forward-chaining materialization in three rounds.
Round~1 applies the ground rules that pair a single asserted triple
with a TBox axiom: \texttt{rdfs:subClassOf},
\texttt{rdfs:subPropertyOf}, \texttt{rdfs:domain}, \texttt{rdfs:range},
\texttt{owl:inverseOf}, \texttt{owl:SymmetricProperty},
\texttt{owl:equivalentClass}, and
\texttt{owl:equivalentProperty}. Round~2 re-applies
\texttt{subPropertyOf} and \texttt{subClassOf} over the triples
produced in round~1 so that subproperty chains propagate through the
inferred graph. Round~3 applies one-hop transitive closure for all
\texttt{owl:TransitiveProperty} instances. Each round is implemented
as a single SPARQL \texttt{INSERT INTO GRAPH} statement wrapped in
\texttt{DEFINE input:inference}, and executes with row-level logging
temporarily disabled (\texttt{log\_enable(2,1)}) so that large
transactions do not exceed Virtuoso's per-transaction log size limit.

This process generates 80{,}907{,}421 materialized triples for a
total of 117{,}950{,}946 triples (a 3.18$\times$
expansion). Materialization completes in approximately 114~minutes
on a server with 1~TB RAM and writes the materialized triples to a
dedicated named graph (\texttt{<https://rubalkhali.science/inferred/>})
so that query answering can be served from materialized facts
without paying a per-query backward-chaining cost. The dominant
inferred predicates are \texttt{sio:isRelatedTo} (22.6M),
\texttt{rdf:type} (14.6M), \texttt{sio:hasAttribute} (11.3M),
\texttt{sio:isAttributeOf} (5.7M), and \texttt{sio:hasMember} (5.6M),
reflecting the propagation of SIO's property hierarchy and inverse
properties across the ABox.

% -----------------------------------------------
\subsection{Schema definition with ShEx}\label{subsec:shex}
% -----------------------------------------------

We used Shape Expressions (ShEx)~\cite{shex} to define structural
schemas for the core entity types. While OWL provides open-world
reasoning and inference capabilities, it cannot readily enforce data
completeness constraints. We developed ShEx schemas to explicitly
validate the graph's structure, ensuring data quality and compliance
before loading the data into the database.

Our ShEx implementation defines shapes for sampling sites, soil
samples, environmental measurements, DNA extracts, amplicon
libraries, sequencing datasets, quality control records, and
taxonomy observations. These shapes guarantee that essential
attributes are invariably present. For instance,
Listing~\ref{lst:shex_site} shows the shape for a sampling site,
which requires an explicit type assertion and a label while
permitting optional coordinates and a free-text description.

\begin{lstlisting}[caption={ShEx shape for Sampling Site validation.},
  label={lst:shex_site}, basicstyle=\ttfamily\scriptsize, frame=single]
<#SamplingSiteShape> {
    rdf:type [rak:0000002] ;  # Must be of type Sampling Site
    rdf:type . * ;                # May have other inferred types
    rdfs:label . ;                # Must have a label
    dcterms:description xsd:string ? ; # Optional description
    geo:asWKT . ? ;               # Optional geographic coordinates
}
\end{lstlisting}

In words, \texttt{<\#SamplingSiteShape>} requires that every
conforming node carry an explicit type assertion to
\texttt{rak:0000002} (Sampling Site) and an \texttt{rdfs:label}; it
permits any number of additional inferred types and optionally
allows a free-text description and a WKT geometry. Records that
omit the mandatory type or label are rejected before they reach the
triple store.

We integrate the ShEx validation into the automated deployment
workflow so that every ABox module is tested against its shape
immediately after generation and before loading into Virtuoso. We
implement the validator with the Apache Jena ShEx library
(version~4.10.0) invoked from Groovy. The workflow terminates with
an explicit error if any instance violates its shape (for example a
sampling site lacking GPS coordinates, or a DNA extract lacking a
concentration value), which functions as a regression test
guaranteeing that modifications to the ETL scripts cannot silently
introduce structural errors.

% ===================================================================
\section{Evaluation}\label{sec:eval}
% ===================================================================

We evaluated the deployed knowledge graph using four complementary
approaches: competency question evaluation via SPARQL queries with
known expected answers; structural validation of every RDF module
against ShEx shapes; logical consistency checking with the ELK
reasoner, run both before and after OWL-Horst materialization; and
usability assessment through a live web portal.

% -----------------------------------------------
\subsection{Competency questions and SPARQL queries}
% -----------------------------------------------

We defined nine competency questions expressed as SPARQL
queries~\cite{sparql} with known expected answers and automated their
execution against the deployed endpoint so that they serve both as
documentation of the intended query workload and as regression tests
during ontology updates. The questions cover retrieval of XRF
measurements, listing sampling sites with coordinates, extraction of
biome assignments via OWL restrictions, retrieval of Light Elements
percentages, DNA extraction records, provenance tracing from soil
samples to sequencing runs, taxonomic composition per run, a
cross-domain join between environmental temperature and taxonomic
abundance, and per-sample sequencing quality control metrics. We
present three representative listings below; the full nine-query
suite is distributed with the source
code~\cite{gigadb_rak,zenodo_rak} and is pre-loaded as
click-to-load examples in the web portal's SPARQL editor.

Listing~\ref{lst:sparql_xrf} retrieves XRF measurements for a specific
site, requiring traversal of the measurement process, quality, and
value entities.

\begin{lstlisting}[caption={SPARQL query to retrieve field XRF
    measurements for Site 10.},
  label={lst:sparql_xrf}, basicstyle=\ttfamily\scriptsize, frame=single]
PREFIX rak: <https://rubalkhali.science/kb/RAK_>
PREFIX sio: <http://semanticscience.org/resource/SIO_>
PREFIX rdfs: <http://www.w3.org/2000/01/rdf-schema#>
SELECT ?processLabel ?analyte ?concentration ?unitLabel
WHERE {
  ?sample a rak:0000021 ;
          rdfs:label ?sampleLabel .
  FILTER(REGEX(?sampleLabel, "Site 10"))
  ?value sio:000232 ?process .
  ?process rdfs:label ?processLabel .
  FILTER(REGEX(?processLabel, "Field", "i"))
  ?sample sio:000008 ?quality .
  ?quality sio:000215 ?value .
  ?value rak:2000012 ?concentration .
  ?value sio:000221 ?unit .
  ?unit rdfs:label ?unitLabel .
  ?quality a ?qualityClass .
  ?qualityClass rdfs:label ?analyte .
  FILTER(?qualityClass != rak:0000029)
}
ORDER BY ?processLabel ?analyte LIMIT 100
\end{lstlisting}

Read in prose, the query first selects every soil sample whose label
matches \texttt{Site~10}, then walks back from each measurement value
to its producing process and forwards from the sample through its
quality and value to the recorded concentration and unit. The final
filter excludes the Light Elements pseudo-analyte
(\texttt{rak:0000029}) so that only individually quantified elements
appear in the result.

Listing~\ref{lst:sparql_taxon} retrieves taxonomic abundance
profiles, joining across bioinformatic workflows, FASTQ datasets,
and taxon observations.

\begin{lstlisting}[caption={SPARQL query to retrieve taxonomic
    composition and abundance for all samples.},
  label={lst:sparql_taxon}, basicstyle=\ttfamily\scriptsize, frame=single]
PREFIX rak: <https://rubalkhali.science/kb/RAK_>
PREFIX sio: <http://semanticscience.org/resource/SIO_>
PREFIX rdfs: <http://www.w3.org/2000/01/rdf-schema#>
SELECT DISTINCT ?run ?taxonLabel ?count ?relativeAbundance
WHERE {
  ?proc a rak:0000071 ;
        sio:000230 ?fastq ;
        sio:000229 ?dsAbs , ?dsRel .
  ?dsAbs a rak:0000074 .
  ?dsRel a rak:0000075 .
  ?fastq rdfs:label ?runLabel .
  BIND(REPLACE(?runLabel, "FASTQ dataset for ", "") AS ?run)
  ?dsAbs sio:000059 ?valAbs .
  ?valAbs rak:2000021 ?count .
  ?valAbs sio:000215 ?qualAbs .
  ?qualAbs sio:000011 ?taxon .
  ?taxon rdfs:label ?taxonLabel .
  ?qualRel sio:000011 ?taxon ;
           a rak:0000072 ;
           sio:000214 ?valRel .
  ?valRel rak:2000020 ?relativeAbundance .
  FILTER(?taxon != ?fastq)
}
ORDER BY ?run DESC(?count) LIMIT 100
\end{lstlisting}

The query identifies every bioinformatic workflow that produces both
an absolute and a relative taxonomic dataset for the same FASTQ
input, then retrieves, for each absolute observation, the read count,
the associated taxon label, and the matching relative abundance from
the parallel relative dataset. Results are grouped by sequencing run
and ordered by descending count.

Listing~\ref{lst:sparql_env_tax} shows a cross-domain query that
joins environmental temperature with taxonomic abundance---a join
that would require at least four separate tables in a relational
database.

\begin{lstlisting}[caption={SPARQL query joining environmental
    temperature and \emph{Pseudomonas} relative abundance.},
  label={lst:sparql_env_tax}, basicstyle=\ttfamily\scriptsize, frame=single]
PREFIX rak: <https://rubalkhali.science/kb/RAK_>
PREFIX sio: <http://semanticscience.org/resource/SIO_>
PREFIX pato: <http://purl.obolibrary.org/obo/PATO_>
PREFIX ncbitaxon: <http://purl.obolibrary.org/obo/NCBITaxon_>
PREFIX rdfs: <http://www.w3.org/2000/01/rdf-schema#>
SELECT ?siteLabel ?temp ?relAbundance
WHERE {
  ?site a rak:0000002 ;
        rdfs:label ?siteLabel .
  ?tempProcess sio:000291 ?site ;
               sio:000229 ?tempVal .
  ?tempVal rak:2000012 ?temp ;
           sio:000215 ?tempQual .
  ?tempQual a pato:0000146 .
  FILTER(?temp > 35.0)
  ?sample sio:000252 ?site .
  ?ext sio:000230 ?sample ; a rak:0000041 ; sio:000229 ?dna .
  ?libPrep sio:000230 ?dna ; a rak:0000062 ; sio:000229 ?lib .
  ?seq sio:000230 ?lib ; a rak:0000064 ; sio:000229 ?fastq .
  ?bioinf a rak:0000071 ;
          sio:000230 ?fastq ;
          sio:000229 ?dsRel .
  ?dsRel a rak:0000075 ;
         sio:000059 ?valRel .
  ?valRel rak:2000020 ?relAbundance ;
          sio:000215 ?qualRel .
  ?qualRel sio:000011 ncbitaxon:286 . # Pseudomonas
}
\end{lstlisting}

The query first selects every sampling site at which the recorded
environmental temperature exceeds 35\,$^{\circ}$C, then for each such
site it walks the full provenance chain---soil sample, DNA extract,
library preparation, sequencing run, FASTQ dataset, bioinformatic
workflow, relative abundance dataset---down to the relative abundance
value of \emph{Pseudomonas} (\texttt{ncbitaxon:286}). Each row of the
result therefore reports a single hot site together with its
corresponding \emph{Pseudomonas} relative abundance.

% -----------------------------------------------
\subsection{ShEx validation results}
% -----------------------------------------------

We validated the structural integrity of every RDF module against
its ShEx shape as described in Section~\ref{subsec:shex}. All eight
shape-annotated modules (\texttt{rubalkhali\_sites.owl},
\texttt{rubalkhali\_samples.owl},
\texttt{rubalkhali\_measurements.owl}, \texttt{rubalkhali\_dna.owl},
\texttt{rubalkhali\_xrf.owl},
\texttt{rubalkhali\_taxonomy\_abox.owl},
\texttt{rubalkhali\_sra.owl}, and \texttt{rubalkhali\_qc.owl}) passed
validation without constraint violations. Because the validation
step is part of the deployment workflow and the workflow aborts on
any violation, any subsequent modification to the ETL scripts that
introduces a structural error is caught before the data reaches
Virtuoso.

% -----------------------------------------------
\subsection{Logical consistency with ELK}
% -----------------------------------------------

We verified the logical consistency of the knowledge graph using
the ELK reasoner~\cite{elk}, an implementation of a
consequence-based reasoning procedure for the OWL~2~EL profile. We
loaded the complete set of OWL modules (base ontology, all ABox
files, and the taxonomy reference graph) into a single OWL ontology
and ran ELK to check for unsatisfiable classes and contradictory
axioms. We ran the reasoner in two configurations: (1)~on the
asserted knowledge graph prior to materialization, and (2)~on the
union of the asserted knowledge graph and the OWL-Horst materialized
graph from Section~\ref{subsec:owl_rdf}.
Table~\ref{tab:elk} reports the results.

In both cases the reasoner confirmed that the ontology is
consistent and reported zero unsatisfiable classes. The
post-materialization run loaded 81.3~million axioms (of which
78.1~million originate from the materialized ABox), discovered 74
additional classes (7{,}202 $\rightarrow$ 7{,}276) from the
materialized type assertions, and required 63~seconds for the
consistency check. These results confirm that the OWL-Horst
entailments introduced no disjointness violations despite adding
5.7~million new individual type assertions, a property that
cannot be guaranteed by inspection alone: new ABox type triples of
the form \(C(a)\) and \(D(a)\) where \(C \sqsubseteq \neg D\) would
render the ontology inconsistent, so the reasoner must be run
after materialization to rule this out.

\begin{table}[h]
\centering
\caption{ELK reasoning performance before and after OWL-Horst
  materialization. All runs used OWL API~4.5.26 and ELK~0.4.3 on a
  server with 1~TB RAM (Java heap 850~GB).}
\label{tab:elk}
\small
\begin{tabular}{lrr}
\toprule
\textbf{Metric} & \textbf{Before} & \textbf{After} \\
\midrule
Modules loaded            & 8             & 10 \\
Total axioms              & 224{,}775     & 81{,}273{,}368 \\
Logical axioms            & 124{,}106     & 17{,}693{,}762 \\
Classes                   & 7{,}202       & 7{,}276 \\
Named individuals         & 30{,}680      & 5{,}712{,}603 \\
\texttt{isConsistent}     & 0.84\,s       & 63.4\,s \\
Class hierarchy           & 0.40\,s       & 0.73\,s \\
Unsatisfiable classes     & 0             & 0 \\
Load time                 & 675\,s        & 7{,}585\,s \\
Peak RSS                  & 203\,GB       & 395\,GB \\
Wall time                 & 688\,s        & 7{,}722\,s \\
\bottomrule
\end{tabular}
\end{table}

We chose ELK over more expressive reasoners such as HermiT~\cite{hermit}
because the combination of TBox and 37-million-triple ABox
renders full OWL~2~DL reasoning computationally infeasible on commodity
hardware. The consequence-based calculus implemented by ELK
produces class-hierarchy results in under two seconds on the asserted
graph, and the post-materialization run confirms the same consistency
result at the cost of a higher load time and increased memory.

% -----------------------------------------------
\subsection{Compute-time metrics}\label{sec:compute_metrics}
% -----------------------------------------------

Table~\ref{tab:compute} summarizes the computational cost of the
key deployment steps. OWL-Horst forward-chaining materialization
runs in approximately 114~minutes and expands the knowledge graph
from 37.0~million asserted triples to 117.9~million
(3.18$\times$). ELK consistency checking over the asserted TBox
takes under two seconds; the post-materialization run that
includes the full materialized ABox requires 63~seconds but
confirms the same result (zero unsatisfiable classes). Competency
question latency on the Virtuoso endpoint ranges from 5~ms for
simple site look-ups (CQ2) to over one second for queries that
traverse inverse properties across the entire DNA extraction
workflow (CQ5).

\begin{table}[h]
\centering
\caption{Compute-time metrics for key deployment and validation
  steps.}
\label{tab:compute}
\small
\begin{tabular}{lr}
\toprule
\textbf{Step} & \textbf{Time} \\
\midrule
RDF module loading (Virtuoso, 9 ABox files)  & 140\,s \\
OWL-Horst materialization (3 rounds)         & 114\,min \\
ELK classification (asserted only)           & 1.6\,s \\
ELK classification (with materialized ABox)  & 64\,s \\
CQ1 -- XRF measurements, Site 10             & 23\,ms \\
CQ2 -- Sampling sites + GPS                  & 5\,ms \\
CQ3 -- Biomes per site                       & 5\,ms \\
CQ4 -- Light Elements percentages            & 8\,ms \\
CQ5 -- DNA extraction records                & 1{,}342\,ms \\
CQ6 -- Provenance chain                      & 255\,ms \\
CQ7 -- Taxonomic composition                 & $>$300\,s \\
\bottomrule
\end{tabular}
\end{table}

The taxonomy composition query (CQ7) exceeds a 5-minute timeout
because it joins absolute and relative abundance datasets across
1.4~million observations and therefore touches a large fraction of
the materialized graph. In practice this query is not used by the
web portal; the portal retrieves taxonomic data per sample rather
than across all runs, which completes in sub-second time. CQ7 is
retained as a stress test for the deployment and as a marker of the
frontier of what can be served interactively from the current
hardware.

% -----------------------------------------------
\subsection{Web portal}
% -----------------------------------------------

We developed a web portal~(\url{https://rubalkhali.science/}) that
provides user-friendly access to the knowledge graph. The portal is
implemented as a React front-end backed by a FastAPI service and
Virtuoso: the backend dynamically generates SPARQL queries to fetch
content directly from the knowledge graph in real time, without
maintaining a separate relational database. Every page rendered by
the portal constitutes a functional test of the underlying data, so
that any data integrity problem surfaces as a visible error. The
portal provides an interactive map of the 70 sampling sites,
per-site dashboards with environmental measurements, geochemistry,
and taxonomic profiles, and a SPARQL query page built on
YASGUI~\cite{yasgui} that provides an autocomplete-aware editor,
prefix management, and a tabbed interface. The nine competency
questions of Section~\ref{sec:eval} are exposed as click-to-load
examples in the SPARQL page so that users can run them against the
live endpoint without typing. The portal additionally offers a
downloads page that serves the individual OWL/RDF modules, the ShEx
shape definitions, and the competency questions as static files
under the same free software licence as the rest of the resource.

% ===================================================================
\section{Discussion}\label{sec:discussion}
% ===================================================================

We revisit the challenges introduced in the introduction, link each
to the design choices described in the preceding sections, and
conclude with a discussion of limitations and sustainability.

% -----------------------------------------------
\subsection{Design rationale}
% -----------------------------------------------

\paragraph{Unified data model.}
The first challenge was to represent heterogeneous observation types
under a single vocabulary. We adopted SIO as the upper-level
framework precisely because it offers a small set of relations
(\texttt{hasInput}/\texttt{hasOutput}, \texttt{hasTarget},
\texttt{hasMember}, \texttt{refersTo}, \texttt{hasValue}) that
generalize across sampling, measurement, and provenance patterns.
Sample collection, XRF analysis, DNA extraction, sequencing, and
taxonomic classification are all instantiations of the same
\textsc{Process--hasInput/hasOutput--Object} idiom; quantitative
results are uniformly represented as
\textsc{Quantity--refersTo--Quality} pairs with explicit units; and
feature tables are represented as
\textsc{Collection--hasMember--Thing} structures. This uniformity is
what allows cross-domain SPARQL queries to traverse from a climate
observation through a sampling site to a taxon abundance without any
per-domain join logic, as demonstrated by the query in
Listing~\ref{lst:sparql_env_tax}. Recent work on simplified
upper-level ontologies, notably SULO~\cite{sulo2025}, takes this
argument further: by strengthening a few fundamental relations and
reducing reliance on specialized predicates, an upper-level ontology
can lower the modelling effort for new domains while preserving
interoperability. Our patterns are compatible with this direction
and could be re-grounded in SULO with limited refactoring.

\paragraph{Taxonomy harmonization.}
The companion data paper~\cite{gigascience_companion} describes our
mapping of operational taxonomic units from multiple
primer/classifier combinations onto NCBI Taxonomy and GTDB. A key
design choice is that we keep both the raw classifier output and the
harmonized identifier in the knowledge graph, so that downstream
consumers can re-use the mapping or substitute their own without
losing the original assertions. This is important for
reproducibility: because classifier outputs are sensitive to
reference database version, a consumer that reruns the classification
against a newer reference may legitimately produce a different
taxon assignment for the same read, and the knowledge graph must
allow both assignments to coexist.

\paragraph{Schema validation at ingest.}
We chose ShEx for ingest validation because the ETL is batch-oriented
and the validator runs in the same Groovy/Jena toolchain that
produces the RDF. SHACL is an obvious alternative and is supported
in-database by several triple stores; a future port of the shapes
from ShEx to SHACL would enable in-database validation at load time
and is a natural extension of the current workflow. We deployed
Virtuoso as the triple store because it is available as free
software under the GNU General Public License, supports OWL-Horst
forward-chaining materialization, and has demonstrated operational
maturity in large-scale linked-data deployments.

\paragraph{Semantic query answering at scale.}
Full OWL~2~DL reasoning over the combined TBox and 37-million-triple
ABox is not computationally feasible with HermiT on commodity
hardware. We therefore split the reasoning task. We use
ELK~\cite{elk} to \emph{validate} the logical consistency of the
knowledge graph (confirming zero unsatisfiable classes both before
and after materialization); we use Virtuoso's OWL-Horst
implementation to \emph{materialize} entailments under pD* semantics
so that query answering can be served from physical triples without
a per-query backward-chaining cost. This split is deliberate: ELK
is sound and complete for classification in OWL~2~EL and therefore
gives us a trustworthy ``no contradiction'' result, but its
consequence-based calculus does not directly produce the kind of
ABox expansion (inverse properties, subproperty chains) that
SPARQL traversal queries benefit from. OWL-Horst, conversely, covers
the rules that matter for query answering---subclass, subproperty,
domain, range, inverse, symmetry, equivalence, transitivity---but is
incomplete with respect to full OWL~2~DL. Applying both reasoners
gives us the best of both: a logically consistent graph whose
materialized form exposes exactly the entailments that the
competency questions need. Two decades of work on scaling OWL
reasoning~\cite{owlhorst,owl2profiles,elk} provide additional
alternative points in the design space, and a more expressive
reasoner could be applied to a sub-module of the TBox in future work.

\paragraph{User-facing interface.}
The web portal addresses the challenge that domain scientists are
typically not fluent in SPARQL. The interactive map and per-site
dashboards cover the most common exploratory tasks; the YASGUI-based
SPARQL editor exposes the underlying graph for power users,
pre-populated with the competency questions as click-to-load
examples, autocomplete for prefixes and predicates, and a tabbed
query interface so that multiple queries can be developed in
parallel.

% -----------------------------------------------
\subsection{Limitations and outstanding challenges}
% -----------------------------------------------

Several challenges remain open. First, although OWL-Horst
materialization plus ELK classification is sufficient for the
queries described in this paper, more complex reasoning---for
example over disjunctions arising from ENVO biome classifications,
or over negative axioms about taxon absence---would require either
a different reasoner or the extraction of a sub-module amenable to
a more expressive calculus such as HermiT. Second, the resource
depends on upstream ontologies (SIO, ChEBI, ENVO, NCBI Taxonomy)
that evolve independently; keeping the deployed graph aligned with
their releases is an ongoing maintenance task and a known
reproducibility risk for any SPARQL query that bottoms out in a
specific external IRI. Third, the 11~GB materialized ABox dump
exceeds the in-memory parsing capacity of the OWL API Turtle
parser on a workstation with less than 128~GB of RAM, so the
post-materialization ELK run currently requires a large-memory
server; splitting the materialized graph into stream-parseable
chunks would relax this requirement. Fourth, CQ7 (the taxonomy
composition query that joins all 1.4~million absolute and relative
abundance observations) still exceeds a five-minute timeout on
commodity hardware; reducing its cost either by query rewriting or
by additional materialization indexing remains future work.

% -----------------------------------------------
\subsection{Sustainability and updates}
% -----------------------------------------------

The dataset is versioned and archived on Zenodo and GigaDB, and the
deployment workflow is reproducible from the public repository. We
intend to extend the resource as additional field expeditions are
processed: new trips will be added as additional ABox modules
without changes to the upper-level patterns. Upstream ontology
updates will be tracked at scheduled intervals; re-materialization
is scripted as the three SPARQL Update statements of
Section~\ref{subsec:owl_rdf} and runs as part of the deployment
workflow, so an upstream release can be absorbed by re-running the
workflow. Where an upstream change would alter the semantics of an
existing competency question, we will record the affected query in
the changelog so that downstream users can audit the impact.

% ===================================================================
\section{FAIR compliance}
% ===================================================================

The design decisions described in the preceding sections collectively
ensure that the resource satisfies the FAIR Data
Principles~\cite{fair}.
\begin{itemize}
  \item \textbf{Findable.} We assigned globally unique IRIs under
    the \texttt{https://rubalkhali.science/kb/} namespace to all
    named individuals, and we registered the dataset with ENA under
    accession PRJEB104209.
  \item \textbf{Accessible.} We serve the knowledge graph via a
    standard SPARQL~1.1 endpoint and a web interface; raw sequencing
    reads are deposited at ENA. Every file is dereferenceable and
    every site and sample IRI resolves to a content-negotiated
    representation.
  \item \textbf{Interoperable.} We employ established shared
    vocabularies (SIO, ENVO, ChEBI, PATO, UO, NCBI Taxonomy) and
    standard formats (OWL~2~DL, RDF/Turtle, SPARQL), and every
    identifier from external ontologies is used as-is so that
    federated queries with other SIO- or OBO-based graphs are
    straightforward.
  \item \textbf{Reusable.} We licensed the data under CC~BY~4.0 and
    described it with rich provenance metadata encoding the
    complete chain from physical sample to digital data product. We
    open-sourced all generation code.
\end{itemize}

% ===================================================================
\section{Conclusion}\label{sec:conclusion}
% ===================================================================

We have presented a formal knowledge graph for the Rub' al-Khali
desert ecosystem that integrates microbial sequencing, soil
geochemistry, environmental observations, and climate records under
a single SIO-based data model. We formalized eight reusable ontology
design patterns---measurement, sampling, experimental protocol,
taxonomic abundance, XRF geochemistry, sampling site with biome,
sequencing workflow, and sequencing quality control---in description
logic and aligned them explicitly with SIO's core relational idioms,
accompanied by Turtle listings and prose verbalizations for every
pattern. We implemented an automated deployment workflow that
generates the modular OWL/RDF files, validates them against ShEx
shapes as a regression test, loads them into Virtuoso, binds the
OWL-Horst rule set, and forward-chains materialization via three
SPARQL Update rounds. The materialization expands the graph from
37.0~million asserted to 117.9~million total triples
(3.18$\times$) in approximately 114~minutes, and the resulting
graph is served through a public SPARQL endpoint and a web portal
with an autocomplete-aware SPARQL editor. We validated the logical
consistency of the full graph with the ELK reasoner both before and
after materialization, confirming zero unsatisfiable classes in both
cases. We evaluated nine competency questions covering per-sample
retrievals, provenance tracing, and cross-domain joins, with
latencies that range from 5~ms for simple site lookups to just over
one second for queries that traverse the DNA extraction workflow.

The resource contributes a reusable blueprint for integrating a
multi-expedition ecological dataset into a Semantic Web
deployment. The ontology design patterns are directly applicable to
other sampling-based environmental studies; the OWL-Horst
materialization workflow is a concrete recipe for any Virtuoso-based
deployment that needs to serve inferred triples without a per-query
reasoning cost; and the combination of ELK for consistency
validation and OWL-Horst for query-time expansion is a practical
template for scaling semantic query answering to tens of millions
of triples on commodity hardware. The underlying dataset is
described in a companion data paper~\cite{gigascience_companion},
and the deployed knowledge graph is available at
\url{https://rubalkhali.science/}.

% ===================================================================
\section{Availability of source code and requirements}\label{sec:availability}
% ===================================================================

All source code for data processing, ontology engineering, RDF
generation, ShEx validation, the SPARQL competency questions, the
React+FastAPI web portal, and the Docker Compose deployment is
publicly available at
\url{https://github.com/bio-ontology-research-group/empty-quarter}
under the MIT licence. The ontology modules, ShEx shapes, competency
questions, and full knowledge graph dump are archived on
Zenodo~\cite{zenodo_rak} and GigaDB~\cite{gigadb_rak} under the
Creative Commons Attribution 4.0 International licence
(CC~BY~4.0). The public SPARQL endpoint is available at
\url{https://rubalkhali.science/sparql}; the interactive web portal is
at \url{https://rubalkhali.science/}.

% ===================================================================
\section*{Declarations}
% ===================================================================

\paragraph{Competing interests.}
The authors declare no competing interests.

\paragraph{Funding.}
Funding statement to be completed.

\paragraph{Author contributions.}
Author contributions to be completed per CRediT.

\paragraph{Acknowledgements.}
Acknowledgements to be completed.

% ===================================================================
\bibliographystyle{plain}
\bibliography{../sn-bibliography}

\end{document}
